\section{Dependence on the barrier}
\label{sec:barrier-dependence}
The preceding logarithmic-rank comparison is not a consequence of
self-concordance alone. Its decisive ingredient is ordered metric
velocity. We first identify a scalar normalization that guarantees this
property, and then show that the same order of centrality cost persists
for explicit coupled barriers with the smallest possible cube parameter.

\subsection{Every fixed normalized scalar barrier has the same sharp factor}
Let $f$ be a $C^3$ standard self-concordant barrier on an open interval
$I=(\ell,b_f)$, with $f''>0$ and an analytic center $x_f\in I$.
Assume that $f'$ maps $(x_f,b_f)$ onto $(0,\infty)$ and that
\begin{equation}
 (f'(x))^2\leq f''(x)\qquad(x_f<x<b_f).
 \label{eq:scalar-normalization}
\end{equation}
The normalization is only required on this positive branch. Define
\[
 \rho_f(x)=\int_{x_f}^x\sqrt{f''(u)}\dd u,\qquad
 h_f(t)=\rho_f((f')^{-1}(e^t)),\qquad v_f=h_f'.
\]
The coordinate change is an isometry for the scalar Hessian metric;
therefore the product $F_f(x)=\sum_{i=1}^r f(x_i)$ has Euclidean
distance in the coordinates $\rho_f(x_i)$. Coordinates below $x_f$
can be moved to $x_f$ while decreasing distance from the center and
improving every positive linear objective. Zero objective coordinates
remain at $x_f$; below, $r$ denotes the number of positive weights.

\begin{theorem}[Normalized scalar comparison]
\label{thm:normalized-scalar}
Under these hypotheses, every positive-objective central subarc satisfies
\[
 L_\CP(\eta_0,\eta_1)\leq
 \Gamma_r d_{F_f}(x(\eta_0),x(\eta_1)).
\]
For each fixed $f$ and $r$, the supremum of the ratio from the analytic
center, over positive objectives and terminal parameters, equals
$\Gamma_r$. Moreover $b_f<\infty$, and for
$0<\eps<(b_f-x_f)\sum_iw_i$,
\begin{equation}
 L_{\CP,F_f}(\eps)\leq C_{\SC}\Gamma_r L_{\opt,F_f}(\eps),
 \qquad C_{\SC}<2.
 \label{eq:general-accuracy}
\end{equation}
Here $L_{\opt,F_f}$ is distance from $(x_f,\ldots,x_f)$ to
$\sum_iw_i(b_f-x_i)\leq\eps$, and the dimension-independent
constant is the explicit attained maximum in
\cref{prop:scalar-envelope}. Numerically, $C_{\SC}\approx1.831856423$.
\end{theorem}
\begin{proof}
At $x=(f')^{-1}(e^t)$, differentiation gives
\begin{equation}
 v_f=\frac{f'}{\sqrt{f''}},\qquad
 v_f'=v_f\left(1-\frac{v_f f'''}{2(f'')^{3/2}}\right)
 \geq v_f(1-v_f)\geq0.
 \label{eq:normalized-profile-ode}
\end{equation}
Thus $v_f$ is nondecreasing and $0<v_f\leq1$. Its limit at
$-\infty$ is zero, by continuity at $x_f$. Its positive limit at
$+\infty$ must be one: a smaller limit would make the derivative
eventually bounded below by a positive number. The negative tail is
integrable because $\int_{-\infty}^t v_f=h_f(t)<\infty$.
Once $v_f\geq1/2$, the differential inequality implies
$(1-v_f)'\leq-(1-v_f)/2$. Consequently
\begin{equation}
 v_f-\mathbf1_{(0,\infty)}\in L^1(\R),\qquad
 h_f(t)=t+\kappa_f+o(1)\quad(t\to\infty).
 \label{eq:general-profile-tails}
\end{equation}
Central metric coordinates are $h_f(s+\log w_i)$. Their velocities
have the order of the weights, so the prefix proof of
\cref{thm:sharp-prefix} applies. For sharpness, set
$h_f(\log w_i)=M(\sqrt i-\sqrt{i-1})$ and terminate at $s=0$.
The integrable step replacement in that proof has error $O_{f,r}(1)$;
distance is $M\Gamma_r$ and length is
$M\Gamma_r^2+O_{f,r}(1)$. This also proves that the supremum is
for each fixed profile, rather than for a profile varying with dimension.

Writing $x(t)=(f')^{-1}(e^t)$ gives
\begin{equation}
 x'(t)=e^{-t}v_f(t)^2,\qquad
 b_f-x(t)=\int_t^\infty e^{-u}v_f(u)^2\dd u.
 \label{eq:general-gap-tail}
\end{equation}
The first formula makes the right endpoint finite; surjectivity of
$f'$ identifies the limit with $b_f$. The second follows by integration.
It remains to prove the accuracy estimate. Put
$p_f(y)=f'(\rho_f^{-1}(y))$ on $y\geq0$. Then
\begin{equation}
 p_f(0)=0,\quad p_f'>0,\quad
 |p_f''|\leq p_f',\quad p_f\leq p_f'.
 \label{eq:general-p-conditions}
\end{equation}
A closest accurate metric vector $y^\star\geq0$ exists: intersect
the closed feasible set with a ball containing one accurate vector,
and use compactness and the finite-coordinate inverse $\rho_f^{-1}$.
The constraint is active, since decreasing a positive coordinate
otherwise reduces the norm. Every coordinate is positive. Indeed,
increasing a zero coordinate by $t$ improves the gap to first order
while increasing squared norm only to second order; compensating by
decreasing a positive coordinate preserves the gap and reduces squared
norm to first order. The constraint gradient is nonzero. Necessary
Lagrange multiplier conditions at this global minimizer therefore give
\begin{equation}
 y_i^\star p_f'(y_i^\star)=\lambda w_i,\qquad \lambda>0.
 \label{eq:general-accuracy-KKT}
\end{equation}
No convexity or uniqueness of this allocation is asserted for a general
$f$. With $k=\log2$, \cref{prop:scalar-envelope} proves
\[
 p_f(y)\leq yp_f'(y)/k\leq p_f(C_{\SC}y).
\]
The center at $\eta=\lambda/k$ thus has metric coordinates between
$y_i^\star$ and $C_{\SC}y_i^\star$. It is accurate, and the first
accurate center is coordinatewise no farther. Apply the endpoint bound.
\end{proof}

More generally, the exact local test for monotonicity is
$2(f'')^2-f'f'''\geq0$ on the positive branch, as follows from
\eqref{eq:normalized-profile-ode}. Monotonicity alone proves the
$\Gamma_r$ upper bound; the normalization additionally supplies the
unit limiting speed, integrable tails, and uniform accuracy comparison.

\begin{proposition}[The full self-concordant scalar class]
\label{prop:nonmonotone}
Over smooth self-concordant interval barriers with an analytic center
and full positive gradient range, the supremum of the
analytic-center-to-central-endpoint ratio in $r$ identical product
coordinates is $\sqrt r$. In the family witnessing sharpness,
the scalar barrier parameter is $\nu_A=\Theta(A^2)$ and $A$ tends
to infinity. Thus this assertion does not concern one fixed normalized
barrier or a dimension-independent parameter budget.
\end{proposition}
\begin{proof}
Every central metric coordinate is increasing, whether or not its
velocity is monotone. Hence
$L\leq\int\norm{\dot y}_1=\norm{y_{\mathrm{end}}}_1
\leq\sqrt r\norm{y_{\mathrm{end}}}_2$.
\Cref{app:nonmonotone-construction} gives a globally smooth barrier
family and disjoint coordinate-velocity windows approaching equality.
\end{proof}

\subsection{The same dyadic separation for every normalized profile}
\begin{corollary}[Fixed-profile finite sequences]
\label{cor:general-scalar-discrete}
Fix $f$ as in \cref{thm:normalized-scalar}. Use exactly the weights
and tolerance \eqref{eq:dyadic-family}, with objective gap
$\sum_iw_i(b_f-x_i)$ and initial point $(x_f,\ldots,x_f)$.
Then
\[
 L_{\opt,F_f}(\eps_r)=\Theta_f(r\sqrt{\log r}),\qquad
 L_{\CP,F_f}(\eps_r)=\Theta_f(r\log r).
\]
For $0<R<1$, arbitrary-label sequences obeying
\eqref{eq:tube-contract} in this metric, with the stated new initial
point and objective gap, satisfy \eqref{eq:growing-tube} whenever
$\delta_r=o(r^{2/3})$. The constants may depend on $f$ and $R$.
In particular, every fixed metric tube, and every fixed Newton-decrement
neighborhood with radius $0\leq\beta<1/2$, has a
$\Theta(\sqrt{\log r})$ overhead in bounded movement.
\end{corollary}
\begin{proof}
Choose $\tau_f$ so that $v_f(t)\geq1/2$ for $t\geq\tau_f$.
By \eqref{eq:general-gap-tail}, each central coordinate has gap at
most $e^{-t}$, and at least $e^{-t}/4$ when $t\geq\tau_f$.
Thus the center at $T$ is accurate and the one at $T-\log8$ is
not, for large $r$: only $O_f(\sqrt r)$ terminal coordinates fail
the lower-tail condition, by \eqref{eq:dyadic-tail}. From
\eqref{eq:general-profile-tails},
$|h_f(t)-t_+|\leq B_f$ for all $t$. Therefore the accurate-center
distance is $\Theta_f(r\sqrt{\log r})$. Every accurate endpoint
has $b_f-x_i\leq\eps_r/w_i$, which implies
\[
 \rho_f(x_i)\geq T-a_i-\log r-O_f(1)
\]
whenever the right side is large. To see this implication directly,
write $x_i=x(t)$ near $b_f$ and combine
$b_f-x(t)\geq e^{-t}/4$ with $h_f(t)=t+O_f(1)$.
The first $\lfloor\sqrt r\rfloor$ coordinates give the matching
lower bound on target distance, as in \cref{thm:dyadic}.

Use the shifted thresholds $a_i+\tau_f$ and cutoff
$a_{J+1}+\tau_f$, where $J=\lfloor r^{2/3}\rfloor$.
The first $j$ velocities exceed $1/2$ after threshold $j$.
The gap spacing and harmonic sum in \eqref{eq:dyadic-tail}
are unchanged by this common shift. They give central length
$\Omega_f(r\log r)$ before the first accurate center; the prefix
bound gives the matching upper bound.
For discrete sequences define the clipped potential $P_f$ by replacing
$0,b$ in the proof of \cref{thm:tube} with
$\tau_f,a_{J+1}+\tau_f$. The exact product metric and velocity
lower bound give each progress increment
$O_f(C_r\sqrt{1+C_r})$, by the same calculation as
\eqref{eq:one-progress-jump}, with $v_0$ replaced by $1/2$.
The start tube gives
$P_f(s_0)\leq2\Gamma_r\delta_r$ by the prefix bound.
Actual accuracy in coordinate one forces
$s_N\geq T-\log r-\delta_r-O_f(1)>a_{J+1}+\tau_f$.
Hence net progress is $\Theta(r\log r)$ even with backtracking labels.
Finally the straight segment in $\rho_f$ coordinates to the center
at $T$ has length $O_f(r\sqrt{\log r})$; chord sampling gives
the unrestricted upper count. \Cref{lem:decrement-tube} applies to
the product barrier and proves the decrement assertion.
\end{proof}
The objective weights and tolerance in this corollary are dyadic with
the same encoding lengths as before. An arbitrary fixed $f$ need not
have an efficient finite
arithmetic oracle; this is a statement about geometric movement.

\subsection{Universal and entropic cube barriers}
The universal barrier of Nesterov and Nemirovski and the entropic
barrier of Bubeck and Eldan are canonical constructions for general
convex bodies. The exact dimension bounds were proved by
Lee and Yue and by Chewi, respectively
\cite{LeeYue2021,BubeckEldan2019,Chewi2023}. Their cube
specializations illustrate the scope of \cref{thm:normalized-scalar}.

\begin{proposition}[Canonical barriers on a cube]
\label{prop:canonical-cube}
For $K=(-1,1)^r$, the universal barrier is $U(x)=\sum_i b(x_i)$
up to an additive constant. The entropic barrier is
$E(x)=\sum_i e(x_i)$, where
\[
 e=a^*,\qquad a(\theta)=\log\frac{2\sinh\theta}{\theta},
 \qquad a(0)=\log2.
\]
Both have exact parameter $r$. For each, the worst endpoint ratio
is exactly $\Gamma_r$, the worst same-accuracy ratio has order
$\sqrt{\log r}$, and \cref{cor:general-scalar-discrete} applies.
\end{proposition}
\begin{proof}
In orthant $\sigma\in\{-1,1\}^r$, the polar of $K-x$ is the
simplex $\sum_i(1-\sigma_ix_i)|z_i|\leq1$, of volume
$[r!\prod_i(1-\sigma_ix_i)]^{-1}$. Summing gives
\[
 \operatorname{vol}((K-x)^\circ)
 =\frac1{r!}\prod_i\left(\frac1{1-x_i}+\frac1{1+x_i}\right)
 =\frac{2^r}{r!}\prod_i(1-x_i^2)^{-1}.
\]
Its logarithm is the claimed universal barrier. Fubini's theorem gives
$\log\int_K e^{\ip{\theta}{z}}\dd z=\sum_i a(\theta_i)$;
separating the supremum in Fenchel conjugacy gives the entropic formula.
These factorization identities are elementary specializations, not
new barrier constructions; product consistency is also explicit in
Chewi's tensorization argument~\cite[Section 4]{Chewi2023}.
The explicit cube log-partition and its coordinatewise gradient also
appear in recent work on linear bandits
\cite[Section 6.1]{LevyValeauAkhavanRebeschini2026}.

For the entropic summand, $a''>0$ and
\[
 x=a'(\theta)=\coth\theta-1/\theta,
 \qquad a''(\theta)=1/\theta^2-1/\sinh^2\theta,
\]
with their continuous values at zero. The map $a'$ increases from
$-1$ to $1$, and $e'(x)=\theta$. Chewi's theorem gives standard
self-concordance with parameter at most one in dimension one.
Directly, the squared local gradient norm is
\[
 \frac{e'(x)^2}{e''(x)}=\theta^2a''(\theta)
 =1-(\theta/\sinh\theta)^2\longrightarrow1
 \quad(\theta\to\infty).
\]
Thus the scalar parameter is exactly one and its positive branch
satisfies \eqref{eq:scalar-normalization}; its center is zero.
The product parameter is exactly $r$ by taking all coordinates to
the positive endpoint. The same facts for $b$ were proved in
\cref{prop:standard-parameters}. Apply the normalized scalar theorem
and its dyadic corollary.
\end{proof}

Alternative interval barriers and their Hessian geometry have substantial
prior literature. Nesterov and Todd~\cite[Section 6.2]{NesterovTodd2002}
use $-\log\cos$ on a rescaled interval, while Papa Quiroz and
Oliveira~\cite{PapaQuirozOliveira2007} construct a different cube
barrier with parameter $3r/2$ and study geometric algorithms.
Neither scalar flattening nor changing the cube barrier is new here.
The conclusion above concerns the sharp central-arc comparison and
the stated finite-sequence separation for the two canonical barriers.

\subsection{Transfer to trace-separable spectral metrics}
\label{subsec:general-spectral-transfer}
The scalar argument extends to the spectral setting without an assumption
that commuting paths are globally shortest. That assertion follows
from the Hessian itself. If $F(x)=\tr f(x)=\sum_i f(\lambda_i(x))$
on a Jordan spectral interval, its Peirce Hessian is
\begin{equation}
 D^2F(x)[u,u]=\sum_i f''(\lambda_i)u_i^2+
 \sum_{i<j}\frac{f'(\lambda_i)-f'(\lambda_j)}{\lambda_i-\lambda_j}
          \norm{u_{ij}}^2.
 \label{eq:general-spectral-Hessian}
\end{equation}
The continuous divided difference is used at coincidences, with the
Peirce normalization in \eqref{eq:jordan-Hessian}. One can obtain
\eqref{eq:general-spectral-Hessian} directly: polynomial Jordan
calculus gives $D\tr f(x)[u]=\ip{f'(x)}{u}$ and makes
$Df'(x)$ act on $V_{ij}$ by the displayed divided difference.
Approximate $f''$ uniformly by polynomials on a compact interval
containing the spectrum and integrate twice. The resulting polynomial
functions and their first two derivatives converge uniformly; their
divided differences do also, by their integral representation.
This proves the formula for $C^2$ functions. It is the trace-separable
case of the classical spectral differential calculus
\cite{LewisSendov2001,Baes2007}.

All off-frame terms are nonnegative. The almost-everywhere spectral
derivative argument in \cref{thm:exact-distance} consequently gives
\[
 \norm{\dot x}_{F,x}^2\geq
 \sum_i f''(\lambda_i)\dot\lambda_i^2
 =\norm{\frac{\dd}{\dd t}\rho_f(\lambda(x))}^2.
\]
Thus center-to-point distance, and distance between coordered endpoints
in one frame, is Euclidean in the transformed spectra. Positive spectral
support objectives have central equations
$f'(\lambda_i)=\eta w_i$ and a common objective frame, so
\cref{thm:normalized-scalar} applies, with active spectral rank.
For the accuracy bound, the Jordan trace inequality aligns an endpoint
with the ordered objective; replacing coordinates below $x_f$ by $x_f$
improves the objective and distance. The scalar minimizer and necessary
conditions \eqref{eq:general-accuracy-KKT} then give the same bound.

For rectangular matrices, assume $f$ is even on $(-b_f,b_f)$,
$x_f=0$, and $f''>0$. Apply the Hermitian dilation argument with
$g=f/2$ to $F(X)=\sum_i f(\sigma_i(X))$; the two copies of each
singular value recover exactly $f''$. The same conclusions hold.
These are assertions about Hessian metrics. Standard self-concordance
of an arbitrary spectral lift is not inferred from scalar
self-concordance. Bounded-move corollaries for a lift require that
property separately; \cref{prop:standard-parameters} supplies it for
the standard spectral barriers treated earlier.
